% !TeX root = Alexander_Alexiev_CV.tex
% Page 2: research, industry, teams, community, and references.

\section{Research Experience}
\entry{Research Assistant in Biomimetic Robotics Laboratory}{Massachusetts Institute of Technology}{Boston, MA}{June 2025 -- present}
\begin{itemize}
  \item Investigating performance gains from quasi-direct-drive (QDD) actuators for sim-to-real transfer in reinforcement learning for dexterous manipulation tasks.
  \item Created a novel method to grasp objects blindly using only hand proprioception via reinforcement learning and model-based optimal control; \textit{See to Reach, Feel to Grasp}.
  \item Developed STM32 firmware for a three-finger tendon-driven hand to decouple joints and enable impedance control at 1,000 Hz over CAN FD.
  \item Developed firmware and a machine-learning training pipeline for soft fingertip sensors that estimate force and contact location at 200 Hz using CAN and LCM.
\end{itemize}

\entry{Undergraduate Researcher in Robot Vision and Learning Lab}{University of Toronto}{Toronto, ON}{June 2022 -- Apr 2025}
\begin{itemize}
  \item Developed an undergraduate thesis technique combining non-rigid point-set registration with a graph neural dynamics model for real-time state estimation and prediction of rope and cloth, robust to occlusion and self-overlap.
  \item Created an RGB-D computer-vision pipeline with a custom CNN to improve a graph-based task-and-motion-planning algorithm.
\end{itemize}

\entry{Robotics Research Co-op in Lab for Translational Engineering}{Harvard Medical School / MIT}{Boston, MA}{June 2023 -- Jan 2025}
\begin{itemize}
  \item Designed an ingestible robot to deliver mRNA vaccines into stomach tissue, including an ultra-low-power microheater array for igniting reactive metallic nanofilm and an actuator that converts the energy into a penetrating jet of mRNA.
  \item Created an ingestible robot with a 10-channel light-sensor array to autonomously diagnose ischemia and gastric bleeding; first author in \textit{Science Robotics}.
  \item Developed an embedded C++ algorithm for an implantable device that detects opioid overdose from live IMU data and autonomously administers naloxone; validated detection in a swine terminal study.
\end{itemize}

\section{Industry Experience}
\entry{Machine Learning Research Intern}{Electronic Arts}{Vancouver, BC (Remote)}{May 2021 -- Sept 2021}
\begin{itemize}
  \item Led development of a computer-vision model for detecting game-engine animation glitches and non-realistic behavior in FIFA and NHL titles.
  \item Designed an autoencoder CNN for anomaly detection on rendered gameplay assets, detecting cloth-animation glitches, missing textures, and irregular human-body configurations.
\end{itemize}

\section{Student Teams}
\cvrow{\textbf{Software Co-Lead}, RoboSoccer Humanoid Robotics Team}{Sept 2022 -- Sept 2023}
\begin{itemize}
  \item Ported bipedal motion planning and control from MATLAB to a custom STM32 implementation, increasing the feedback loop to 1,000 Hz; designed a decision-tree-based multi-agent controller for four humanoid robots at RoboCup 2022.
\end{itemize}
\cvrow{\textbf{Embedded Systems and Controls Team}, University of Toronto Formula Racing}{Sept 2020 -- Sept 2022}
\begin{itemize}
  \item Wrote ATmega 2560 firmware for I2C and CAN sensor data and vehicle subsystems including throttle and ABS; designed a cell-level fusing system for a custom 600 V Li-ion battery pack.
\end{itemize}

\section{Community Involvement}
\cvrow{\textbf{FIRST Robotics Team 7558}}{Sept 2018 -- present}
\begin{itemize}
  \item Teach modern robotics methods and tools to high-school students. Currently developing a C++ localization library using NVIDIA Isaac Sim, ROS 2, multiple cameras, AprilTags, an extended Kalman filter, and factor graphs.
\end{itemize}

\section{References}
\begin{tabularx}{\textwidth}{@{}Xr@{}}
  \textbf{Dr. Sangbae Kim}, Professor of Mechanical Engineering & Massachusetts Institute of Technology\\
  \textbf{Dr. Giovanni Traverso}, Professor of Mechanical Engineering & Massachusetts Institute of Technology\\
  \textbf{Dr. Florian Shkurti}, Professor of Computer Science & University of Toronto
\end{tabularx}
